\documentclass[../main.tex]{subfiles}
\begin{document}

\chapter{System Architecture}
\label{ch:architecture}

\section{Design Philosophy}
\label{sec:design-philosophy}

The architecture of AMATAS departs from the prevailing paradigm of network intrusion detection as a single-model classification problem. Rather than training one classifier to distinguish benign from malicious traffic, AMATAS distributes the analytical burden across specialised reasoning agents, each attending to a distinct dimension of the evidence space, whose outputs are then reconciled by a consensus mechanism that explicitly models disagreement. This design is based on the hypothesis that distributed deliberation among diverse models improves classification accuracy on ambiguous data compared to the direct output of a single classifier.

The conceptual foundation for this approach draws on \textcite{Minsky1986} Society of Mind framework, which proposes that intelligence emerges not from any individual reasoning process but from the interaction and negotiation of many smaller specialised agents. Applied to network intrusion detection, this principle motivates a division of analytical labour along the natural fault lines of the problem domain: protocol-layer anomalies, statistical distributional shifts, behavioural pattern matching against known attack signatures, and temporal correlations across flows from the same source. Each of these dimensions demands different domain knowledge, different feature attention, and different reasoning heuristics. Consolidating these tasks into a single prompt increases context complexity, which can degrade the model's performance and accuracy on specialised flow analysis tasks. Separating them into dedicated agents, each with a system prompt tuned for its particular analytical mode, allows each agent to reason with the focused rigour that its subdomain warrants.

Empirical support for multi-agent deliberation in LLM contexts comes from \textcite{Du2023}, who demonstrated that when multiple LLM instances debate a question across several rounds, the accuracy of the final answer significantly exceeds that of any individual instance. \textcite{Chan2023} extended this finding by showing that role diversity among debating agents, rather than merely numerical redundancy, is the critical variable: agents assigned distinct analytical perspectives outperform agents assigned identical roles. These findings translate directly to the AMATAS design, where the four specialist agents are not interchangeable replicas but genuinely different analytical characters operating on different subsets of the feature space. Both Du et al.\ and Chan et al.\ note that multi-agent debate can converge on confidently incorrect answers when agents share initial misconceptions; AMATAS mitigates this through the Devil's Advocate mechanism (Section~\ref{sec:devils-advocate}), which provides a structurally independent counter-voice not subject to the same convergence dynamics.

The adversarial dimension of the architecture is grounded in a distinct body of literature concerned with the failure modes of consensus under social pressure. When agents reason sequentially and share intermediate conclusions, later agents tend to anchor on earlier ones, producing a cascade of agreement that amplifies initial errors rather than correcting them \autocite{Sunstein2002}. Sunstein qualifies this as a ``statistical regularity'' rather than an inevitability, noting that views formed after extensive deliberation are least susceptible to polarisation. The AMATAS design treats polarisation as a sufficiently frequent risk to warrant structural mitigation rather than relying on specialist diversity alone to prevent it. AMATAS addresses this through the Devil's Advocate agent, a dedicated adversarial voice that receives all specialist findings simultaneously and is explicitly instructed to construct the strongest possible benign interpretation of the evidence, regardless of the agent's own internal assessment. \textcite{Yin2024} showed that LLM-powered devil's advocate roles measurably improve group decision quality and reduce inappropriate reliance on confident but incorrect assessments. In the NIDS context, where the cost of a false positive (interrupting a legitimate business process, eroding analyst trust in the system, generating compliance risk) is substantial, this adversarial check provides a principled structural mechanism for false positive suppression rather than relying solely on post-hoc threshold tuning.

The result is an architecture with two orthogonal quality mechanisms: specialist depth, which increases the likelihood of detecting genuine attacks, and adversarial validation, which mitigates false positive classifications. The Orchestrator agent at the apex of the pipeline then faces a well-posed synthesis problem: it receives evidence and counter-evidence from parties with explicit, different mandates, and is required to adjudicate between them according to explicit consensus rules. This architecture functions as an evidence-weighing mechanism rather than a standard ensemble vote, explicitly preserving the reasoning chain at each stage.

\begin{figure}[H]
\centering
\resizebox{0.95\textwidth}{!}{%
\begin{tikzpicture}[
    node distance=8mm and 10mm,
    every node/.style={font=\normalsize},
    box/.style={draw, rounded corners=3pt, minimum height=11mm, text centered, fill=#1!12, draw=#1!60, minimum width=30mm},
    box/.default=black,
    arr/.style={-{Stealth[length=2.5mm]}, thick},
]

% Input
\node[box=black, minimum width=48mm] (input) {Raw NetFlow Record};

% Tier 1
\node[box=Maroon, below=10mm of input, minimum width=62mm] (rf) {\textbf{Tier 1:} Random Forest ($p < 0.15$)};

% Tier 1 outputs
\node[box=ForestGreen, below left=12mm and 15mm of rf, minimum width=36mm] (autoben) {Auto-BENIGN};
\node[below=3mm of autoben, font=\small\itshape, text=gray] {${\sim}95\%$ of flows};

% Phase 1 agents — generous spacing
\node[box=RoyalBlue, below right=20mm and -22mm of rf] (proto) {Protocol};
\node[box=RoyalBlue, right=6mm of proto] (stat) {Statistical};
\node[box=RoyalBlue, right=6mm of stat] (behav) {Behavioural};
\node[box=RoyalBlue, right=6mm of behav] (temp) {Temporal};

% Phase 1 label
\node[above=2mm of proto.north west, anchor=south west, font=\small\bfseries, text=RoyalBlue] {Phase 1 (parallel)};

% Phase 2
\node[box=RedOrange, below=16mm of $(stat.south)!0.5!(behav.south)$, minimum width=58mm] (da) {\textbf{Devil's Advocate} (30\% weight)};
\node[above=2mm of da.north west, anchor=south west, font=\small\bfseries, text=RedOrange] {Phase 2};

% Phase 3
\node[box=Purple, below=16mm of da, minimum width=58mm] (orch) {\textbf{Orchestrator} (consensus)};
\node[above=2mm of orch.north west, anchor=south west, font=\small\bfseries, text=Purple] {Phase 3};

% Output
\node[box=black, below=10mm of orch, minimum width=62mm] (output) {Verdict + Reasoning Chain};

% Arrows
\draw[arr] (input) -- (rf);
\draw[arr] (rf.south) -- ++(0,-5mm) -| (autoben.north);
\draw[arr] (rf.south) -- ++(0,-5mm) -| ($(proto.north)!0.5!(temp.north)$) coordinate (mid1);
\draw[arr] (mid1) -- (proto.north);
\draw[arr] (mid1) -- (stat.north);
\draw[arr] (mid1) -- (behav.north);
\draw[arr] (mid1) -- (temp.north);

% Label on right branch — positioned above the agents
\node[above=2mm of temp.north east, anchor=south west, font=\small\itshape, text=gray] {${\sim}5\%$ suspicious};

\draw[arr] (proto.south) -- ++(0,-5mm) -| (da.north west);
\draw[arr] (stat.south) -- ++(0,-5mm) -| ([xshift=-6mm]da.north);
\draw[arr] (behav.south) -- ++(0,-5mm) -| ([xshift=6mm]da.north);
\draw[arr] (temp.south) -- ++(0,-5mm) -| (da.north east);

\draw[arr] (da) -- (orch);
\draw[arr] (orch) -- (output);

\end{tikzpicture}%
}
\caption{AMATAS two-tier architecture. Tier~1 filters ${\sim}95\%$ of flows as auto-benign; remaining flows pass through the three-phase LLM pipeline.}
\label{fig:architecture}
\end{figure}

\section{Tier 1: Random Forest Pre-Filter}
\label{sec:tier1-rf}

The economic case for LLM-based intrusion detection depends entirely on solving a cost problem that pure LLM pipelines cannot solve alone. At the pricing applicable during Stage~1 experiments (approximately US\$0.028 per flow processed through the full six-agent pipeline), analysing the complete CICIDS2018 dataset of twenty million flows would cost in excess of half a million dollars. This computational cost limits production viability. While LLM reasoning is most beneficial for ambiguous cases, applying it uniformly to all traffic, which is predominantly benign, is computationally and financially inefficient.

The Tier~1 pre-filter resolves this tension by delegating the easy cases to a cheap, fast machine learning classifier and reserving the expensive LLM pipeline for flows that the classifier cannot confidently resolve. The classifier of choice is a Random Forest, a well-understood ensemble method with strong empirical performance on structured tabular data, fast inference, and a natural probabilistic output in the form of class probability estimates derived from voting across an ensemble of decision trees \autocite{Breiman2001}. A Random Forest can process hundreds of thousands of flows per second on commodity hardware, making it suitable for the high-throughput pre-screening role that Tier~1 requires.

\subsection{Feature Set and Training Data}
\label{sec:feature-set-training}

The classifier is trained on twelve numeric features extracted from the CICIDS2018 NetFlow v3 dataset:

\begin{quote}
\small\raggedright
\texttt{L4\_SRC\_PORT}, \texttt{L4\_DST\_PORT}, \texttt{PROTOCOL},
\texttt{FLOW\_DURATION\_MILLI\-SECONDS},
\texttt{SRC\_TO\_DST\_IAT\_MAX}, \texttt{DST\_TO\_SRC\_IAT\_MAX},
\texttt{IN\_BYTES}, \texttt{OUT\_BYTES}, \texttt{IN\_PKTS},
\texttt{OUT\_PKTS}, \texttt{TCP\_FLAGS}, \texttt{DNS\_QUERY\_ID}
\end{quote}

\noindent This feature set was chosen through a combination of variance analysis and domain considerations. Features with near-zero variance across the dataset provide no discriminative signal and were excluded. The \texttt{FLOW\_DIRECTION} field, present in other NetFlow variants, is unavailable in NetFlow v3 and could not be included. \texttt{DNS\_QUERY\_ID} is retained despite being zero for the majority of non-DNS flows, as its presence provides discriminative information for DNS-tunnelling and related attack patterns.

The training data is drawn from the development partition of the CICIDS2018 dataset, comprising 7.04 million flows, which was further stratified into an 80 per cent training set of 5,632,348 flows and a 20 per cent evaluation holdout of 1,408,087 flows. This internal partitioning was necessitated by a data integrity concern identified during the Stage~1 evaluation phase. Seven of the fourteen attack types in the evaluation programme are present in the development CSV, meaning that if evaluation batches were sampled from the same file on which the Random Forest was trained, the classifier's filtering accuracy for those attack types would be inflated by its prior exposure to the same flow distributions. To eliminate this within-split overlap, the RF is trained exclusively on the 80 per cent partition, and evaluation batches for development-resident attack types are sampled exclusively from the 20 per cent holdout. The remaining seven attack types, sourced from the validation and test splits that the RF never encountered during training, require no modification. This partitioning ensures that the Tier~1 evaluation reflects genuine generalisation rather than dataset memorisation.

The Random Forest is configured with one hundred estimators, a maximum tree depth of twenty, a minimum of ten samples per leaf, balanced class weights to account for the strong class imbalance in the training data (approximately eighty-seven per cent benign flows), and a fixed random seed of 42 for reproducibility. Training on the 5.6 million row training partition completes in several hours on a standard multi-core machine.

\subsection{Threshold Selection}
\label{sec:threshold-selection}

The Tier~1 classifier produces, for each flow, an estimated probability that the flow is malicious, derived by averaging the proportion of attack-class votes across all one hundred trees. This probability is compared against a fixed threshold of 0.15. Flows where the estimated attack probability falls below 0.15 are classified as auto-BENIGN and routed away from the LLM pipeline without further analysis. Flows at or above the threshold are forwarded to Tier~2.

The threshold of 0.15 was selected to satisfy an asymmetric optimisation objective: the Tier~1 filter must achieve one hundred per cent recall on attack flows (no attack may be discarded before reaching the LLM pipeline) whilst maximising the proportion of benign flows filtered out. This is the recall-priority constraint appropriate for a security application where false negatives carry the highest cost. At the 0.15 threshold, the RF achieves one hundred per cent recall on the evaluation holdout, with the minimum estimated attack probability across all true attack flows of 0.64. The gap between the threshold (0.15) and the minimum true-attack probability (0.64) provides a substantial safety margin: even if the RF were to slightly miscalibrate its probability estimates on novel flow distributions, attacks would still fall well above the filter threshold. Approximately ninety-five per cent of benign flows are assigned probabilities below 0.15, yielding a filtering rate that reduces LLM API costs by the same proportion.

It bears emphasis that the Tier~1 filter carries a deliberate asymmetry. False negatives at this stage (benign flows forwarded to the LLM pipeline) carry a cost measured in dollars and milliseconds. False positives at this stage (attack flows filtered before LLM analysis) carry a cost measured in missed detections, a fundamentally different and unacceptable failure mode. The threshold of 0.15, far below the intuitive operating point of 0.50, reflects this asymmetry explicitly and structurally.

\section{Tier 2: Specialist Agents}
\label{sec:tier2-specialists}

Flows that survive the Tier~1 filter proceed to the six-agent LLM pipeline, which is organised into three sequential phases. The first phase comprises four specialist agents that execute in parallel, each receiving the same flow record and returning an independent structured analysis. This parallel architecture is both practically efficient (all four API calls are dispatched simultaneously, with wall-clock latency determined by the slowest single call rather than their sum) and analytically important, as it ensures that each specialist reasons from the raw flow data without being influenced by any other agent's intermediate conclusions.

All agents in the pipeline are implemented using GPT-4o, selected after empirical comparison with GPT-4o-mini and Claude Sonnet on a one-hundred-flow validation batch. GPT-4o-mini was rejected after demonstrating a systematic classification bias that produced one hundred per cent false positive rates across the validation set, classifying all flows as malicious regardless of evidence. Claude Sonnet was rejected on operational grounds: its rate limits were insufficient to process one-thousand-flow evaluation batches without prohibitive delays. GPT-4o provided the best trade-off between detection accuracy, false positive control, and throughput at its operating price of US\$2.50 per million input tokens and US\$10.00 per million output tokens.

Each specialist agent is structured around a system prompt that defines the agent's analytical mandate, the specific features it should prioritise, the pattern vocabulary it should employ, and the JSON schema it must produce. The JSON output requirement is explicit in every system prompt and is enforced at the application layer through parsing logic with retry on malformed responses. Every specialist output contains a verdict (one of BENIGN, SUSPICIOUS, or MALICIOUS), a confidence score on the unit interval, a predicted attack type or null, a list of key findings expressed as evidence strings, a full reasoning narrative, and token consumption and cost data for cost accounting.

\textbf{Protocol Agent.} The Protocol Agent's analytical mandate centres on the consistency and validity of the protocol-layer features in the flow record. It examines whether the source and destination ports conform to expected service assignments, whether TCP flag combinations are syntactically valid and semantically coherent for the connection phase they imply, and whether the observed protocol number matches the port and behaviour patterns associated with that protocol. A connection from an ephemeral source port to destination port 21 bearing SYN flags repeated dozens of times within seconds presents a very different protocol picture from a single long-lived FTP control session, and the Protocol Agent is specifically prompted to reason about these distinctions. Port anomalies (traffic to well-known service ports during off-hours, use of non-standard ports for common protocols, combinations of ports associated with specific attack toolkits) are weighted heavily in the Protocol Agent's verdict.

\textbf{Statistical Agent.} The Statistical Agent focuses on the volumetric and timing characteristics of the flow: byte counts in both directions, packet counts, inter-arrival time extremes, and flow duration. Its analytical frame is distributional rather than rule-based. Byte volumes that fall several orders of magnitude above or below typical traffic for the observed protocol, packet counts that imply either atomic single-packet probes or sustained floods, inter-arrival times that collapse to near-zero (indicative of automated generation) or stretch to extreme values (indicative of slow-rate attacks); these distributional aberrations are the primary signal the Statistical Agent is trained to detect. The agent is instructed to reason about what typical traffic of the observed type would look like, and to articulate why the observed values deviate from that expectation, ensuring that its reasoning is anchored in domain knowledge rather than purely in numerical thresholds.

\textbf{Behavioural Agent.} The Behavioural Agent operates at the pattern-matching layer, correlating the full feature profile of the flow against a structured vocabulary of known attack signatures derived from intrusion detection research and mapped to the MITRE ATT\&CK framework \autocite{MITRE2023}. Its prompt enumerates the characteristic behavioural signatures of each attack category present in the CICIDS2018 dataset: the repetitive authentication failure pattern of brute-force attacks, the zero-response-body high-rate request pattern of HTTP floods, the asymmetric byte-volume signature of SYN floods, the slow-rate low-bandwidth pattern of Slowloris attacks. When the Behavioural Agent identifies a match, it returns the corresponding MITRE ATT\&CK technique identifiers alongside its verdict, enabling downstream mapping of detections to the standardised adversary behaviour ontology that security teams use for threat reporting and incident response.

\textbf{Temporal Agent.} The Temporal Agent is architecturally distinct from the other three specialists in that it does not analyse the target flow in isolation. Before invoking the Temporal Agent, the pipeline queries the current batch for all other flows sharing the same source IP address, extracts a compact summary of each co-IP flow, and injects this context into the agent's input. The Temporal Agent therefore reasons about the target flow against the backdrop of the source IP's complete observed behaviour within the batch, identifying patterns (burst activity, sequential escalation from scanning to exploitation, repetitive connection signatures indicative of automated tools, fan-out patterns suggesting lateral movement) that are invisible at the level of individual flows.

This contextual enrichment comes at a cost. The Temporal Agent's input tokens are substantially larger than those of the other specialists, since the co-IP flow summaries can add hundreds or thousands of tokens depending on how active the source IP is within the batch. This makes the Temporal Agent the most expensive single component in the pipeline, accounting for approximately thirty per cent of total per-flow LLM expenditure. Despite this cost, the Temporal Agent is retained as a non-optional pipeline stage because the attack patterns it is uniquely positioned to detect (particularly brute-force sequences and systematic port sweeps) are among the most practically significant in the CICIDS2018 dataset, and no other agent in the pipeline has access to cross-flow context.

\section{The Devil's Advocate Mechanism}
\label{sec:devils-advocate}

Following the completion of the four parallel specialist analyses, the pipeline enters its second phase with the invocation of the Devil's Advocate agent. This agent occupies an unusual position in the architecture: unlike the specialists, whose mandate is to characterise the flow as accurately as possible from their particular analytical angle, the Devil's Advocate's mandate is explicitly adversarial and one-directional. It is instructed to argue for the benign interpretation of the flow regardless of the weight of specialist evidence to the contrary.

The theoretical justification for this asymmetric mandate draws on the literature on group decision-making under uncertainty. When a group of expert agents, human or artificial, is presented with a question and reaches a preliminary consensus, subsequent deliberation tends to reinforce rather than challenge that consensus, a phenomenon known as group polarisation \autocite{Sunstein2002}. In the AMATAS context, if three of four specialists return MALICIOUS verdicts, a naive synthesis agent that simply weighted these verdicts would be strongly predisposed toward a malicious classification before the evidence had been properly stress-tested. The Devil's Advocate agent mitigates group polarisation by systematically generating counter-arguments; it receives all four specialist outputs simultaneously and computes alternative benign explanations for each finding that the specialists cite as suspicious.

\textcite{Yin2024} provide the most directly relevant empirical grounding for this design, demonstrating in a series of human-AI decision experiments that an LLM playing an explicit devil's advocate role measurably improved the quality of group decisions, reduced inappropriate reliance on confident AI assessments, and increased the probability that correct but minority-held positions received adequate consideration before a final verdict. The effectiveness of devil's advocate interventions may vary with task difficulty, a possibility the Stage~1 ablation study (Section~\ref{sec:agent-ablation}) directly examines, revealing that the DA's contribution is indeed non-uniform across attack categories of differing ambiguity. Their findings apply with particular force in the false positive problem of intrusion detection: the cost of classifying a legitimate business flow as malicious (blocking or alerting on an innocent connection, potentially disrupting operations) is substantial and asymmetric relative to the cost of a missed detection in many operational contexts.

The Devil's Advocate agent generates four structured outputs. The \texttt{benign\_argument} field contains a full-text argument for the benign interpretation of the flow, engaging specifically with the evidence cited by the specialists. The \texttt{confidence\_benign} field encodes a normalised score from zero to one representing the strength of the benign case as the Devil's Advocate assesses it. The \texttt{alternative\_explanations} list provides a point-by-point rebuttal of each specialist finding, offering a plausible innocent explanation for each cited anomaly. The \texttt{strongest\_benign\_indicator} field identifies the single most compelling reason that the flow might be benign, which the Orchestrator uses as a key input to its weighting of the Devil's Advocate voice.

The weight assigned to the Devil's Advocate in the final consensus is thirty per cent. This figure was determined empirically through a phase of weight calibration experiments. At fifty per cent Devil's Advocate weight, the adversarial counter-arguments were sufficiently influential to override clear specialist consensus, producing a systematic degradation in recall as genuinely malicious flows were reclassified as benign under the pressure of well-constructed alternative explanations. At thirty per cent, the Devil's Advocate retains sufficient influence to prevent false positives when specialist consensus is weak or founded on borderline evidence, whilst strong agreement among three or four specialists is robust enough to withstand even a confident benign argument. This calibration is empirically validated across the Phase~3 experimental series and held constant across all Stage~1 production experiments.

\section{Orchestrator and Consensus}
\label{sec:orchestrator-consensus}

The Orchestrator agent, the final stage of the pipeline, receives the structured outputs of all five preceding agents (four specialists and the Devil's Advocate) and synthesises them into a single final verdict accompanied by a full reasoning narrative. Its role is not to add new analysis of the flow data itself, but to perform the adjudicative function of weighing competing evidence and resolving conflicts according to explicit, auditable rules.

The consensus rules encoded in the Orchestrator's system prompt implement a threshold-based voting scheme with Devil's Advocate counterweighting. When all four specialists return MALICIOUS verdicts and the Devil's Advocate's \texttt{confidence\_benign} is low, the Orchestrator returns MALICIOUS with high confidence. When three of four specialists return MALICIOUS and the Devil's Advocate argument is moderate, the result remains MALICIOUS but at reduced confidence. When three of four specialists return MALICIOUS but the Devil's Advocate raises a strong alternative explanation, the Orchestrator may downgrade the verdict to SUSPICIOUS. When only two specialists return a MALICIOUS verdict, the Orchestrator assigns a SUSPICIOUS classification unless the specialists' confidence scores exceed a defined threshold. A single MALICIOUS specialist vote is insufficient for a malicious classification in most circumstances, and unanimous BENIGN verdicts produce a BENIGN classification without LLM ambiguity. The SUSPICIOUS verdict category serves as an important intermediate state, capturing flows where the evidence is genuinely ambiguous; in operational deployment, SUSPICIOUS flows would be escalated for human analyst review rather than automatically blocked or ignored.

The Orchestrator also computes a \texttt{consensus\_score}, a normalised measure of specialist agreement, by examining the proportion of specialists whose verdicts align with the final classification. This score provides a machine-readable signal of how confident the ensemble is in its conclusion, separate from the verdict-level confidence score, which reflects the Orchestrator's calibrated belief in the correctness of the verdict given all available evidence. The \texttt{mitre\_techniques} field in the Orchestrator output consolidates all MITRE ATT\&CK technique identifiers cited by the Behavioural Agent and any other specialists that reference the framework, producing a unified list of applicable techniques that security operations teams can use to contextualise the detection within their threat intelligence workflows.

Critically, the Orchestrator's reasoning field is not a summary of the preceding analyses but a synthesis that cites specific findings from specific agents, articulates the evidential basis for the verdict, identifies which agents agreed and which dissented, and explains why dissenting views were or were not dispositive. This full reasoning chain, preserved in the output record for every flow processed through the pipeline, is the primary deliverable from the explainability perspective that motivates the system's design. Unlike a traditional classifier, which produces a numerical probability that offers no foothold for analytical scrutiny, the Orchestrator's output can be read by a security analyst, challenged, corrected, and used as the basis for a documented investigation narrative, providing a traceable decision record that aligns with modern security governance requirements.

\section{Cost Architecture}
\label{sec:cost-architecture}

The economic structure of the Tier~2 pipeline reflects the design decisions described above, with the Temporal Agent's contextual enrichment producing a cost distribution that departs markedly from the even split one might expect across six nominally equivalent agents. Empirical cost data collected across the Stage~1 evaluation experiments reveals that the Temporal Agent accounts for approximately thirty per cent of the total per-flow LLM expenditure, consuming substantially more input tokens than any other agent due to the co-IP flow context injected into its prompt.

The Orchestrator agent is the second most expensive component, at approximately 21.3 per cent of total per-flow cost. This elevated cost reflects the Orchestrator's receipt of all five preceding analyses, each of which may contain several hundred tokens of reasoning text, making its input context the largest of any single agent in the pipeline. The Devil's Advocate follows at approximately 18.5 per cent, since it similarly receives all four specialist outputs. The three remaining specialists (Behavioural, Statistical, and Protocol) operate on the flow data alone and share the remaining thirty per cent approximately evenly, with the Behavioural Agent at 10.4 per cent, the Statistical Agent at 10.1 per cent, and the Protocol Agent at 9.7 per cent. The Behavioural Agent's marginally higher cost reflects the richer attack signature vocabulary in its system prompt, which generates longer input contexts even without additional flow context.

The practical implication of this cost distribution is that the Temporal Agent is the highest-return target for cost optimisation in future iterations of the system. If the contextual enrichment mechanism could be made more selective (providing co-IP context only when the source IP has an unusually high flow count, or compressing the co-IP summaries more aggressively for large groups), the overall per-flow cost could be reduced without fundamentally altering the pipeline's analytical capabilities. Conversely, the three lightweight specialists contribute disproportionately to analytical diversity relative to their cost share, suggesting that their inclusion is strongly justified on cost-effectiveness grounds even if reducing the pipeline to fewer agents were otherwise desirable.

The combined effect of Tier~1 filtering and the Tier~2 cost structure is a system where the effective cost per raw network flow (across all flows entering the system, including those filtered at Tier~1) is approximately five per cent of the per-flow LLM cost. At the observed Stage~1 processing cost of approximately US\$0.028 per flow forwarded to Tier~2, the blended cost across all flows is approximately US\$0.0014 per raw flow, a reduction of approximately ninety-five per cent relative to applying the full LLM pipeline uniformly. This blended cost remains non-trivial at enterprise scale, but it represents the threshold of practical viability in a way that unfiltered LLM analysis does not.

\section{Temporal Clustering Module}
\label{sec:temporal-clustering}

The two-tier architecture described above constitutes the AMATAS v2 system, which is the primary subject of the Stage~1 evaluation experiments. The v3 system, currently in the design phase, extends this architecture with a temporal clustering module inserted between Tier~1 filtering and the Tier~2 LLM pipeline. The motivation for this extension is a hypothesis grounded in the performance profile observed across Stage~1 attack categories.

Certain attack types that operate by generating individually ambiguous flows (most notably DoS-Hulk and Infiltration) present a fundamental challenge to flow-level LLM analysis. Each individual flow, examined in isolation or even in the context of other flows from the same source IP within a single batch, may present no clearly anomalous features. The malicious character of these attacks emerges only when flows are examined in aggregate across time: a sustained pattern of high-rate HTTP requests to a single server, each individually plausible but collectively characteristic of a volumetric attack, or a sequence of reconnaissance, credential testing, and lateral movement flows that spans hours or days. The Temporal Agent's current implementation, which groups flows by source IP within a batch, provides partial temporal context but is limited by the batch boundaries imposed by the experimental evaluation framework.

The v3 Temporal Clustering Module addresses this limitation through two complementary clustering modes. In IP-level clustering mode, flows are grouped by source IP address across a configurable time window (tentatively set at thirty minutes for the initial design, with the window length as a hyperparameter to be tuned on the validation set), and the resulting cluster is treated as the unit of analysis rather than the individual flow. The LLM pipeline receives the full cluster as context, enabling the agents to reason about the collective behaviour of a source IP across extended time periods and across the batch boundaries that constrain the current system. In time-windowed clustering mode, flows are grouped by temporal proximity within a sliding window regardless of source IP, capturing patterns of simultaneous activity from multiple sources that may indicate coordinated attacks such as distributed denial of service.

A design consideration specific to the CICIDS2018 dataset is the role of DNS flows in clustering. DNS flows are extremely frequent in enterprise network traffic and are almost always benign, but they share source IPs with application-layer attack flows in ways that inflate cluster sizes and add noise to the temporal context provided to the LLM agents. The v3 design therefore incorporates a DNS threshold override: flows where DNS\_QUERY\_ID is non-zero (indicating a DNS transaction) are exempt from the clustering mechanism and are routed directly to the standard batch-level Temporal Agent processing unless they themselves exhibit anomalous characteristics. This prevents DNS-heavy source IPs from generating prohibitively large cluster contexts whilst preserving the ability to detect DNS-based attacks such as tunnelling when the DNS traffic itself is anomalous.

The central hypothesis motivating the v3 design is that temporal cluster context will improve recall on attack types where individual flow analysis reaches a ceiling. If a source IP has generated three hundred HTTP requests to the same destination server within a five-minute window, each individually matching the traffic profile of legitimate web browsing, the cluster-level LLM context will reveal the volumetric and temporal pattern that individual flow analysis cannot see. This hypothesis will be tested in the v3 evaluation phase using the held-out test partition of the CICIDS2018 dataset, with comparison against v2 results to isolate the contribution of temporal clustering to detection performance. The expected outcome is an improvement in recall for DoS-Hulk, DoS-Slowloris, and Infiltration, attack categories where the v2 evaluation suggests a fundamental ceiling imposed by the absence of temporal cluster context, without significant degradation in false positive rates, since the cluster context provides the Devil's Advocate with richer benign evidence for flows in genuinely normal clusters.

\onlystandalone{\printbibliography}
\end{document}
